\chapter{Observabilidad, confiabilidad y costo}
\label{chap:observabilidad}

\begin{objetivos}
El estudiante podrá instrumentar señales correlacionadas, formular indicadores y objetivos de servicio, usar presupuestos de error y validar hipótesis operativas con OpenTelemetry, Prometheus, Grafana y Jaeger.
\end{objetivos}

\section{Observar para explicar}

Monitorear responde preguntas previstas; observar permite investigar estados no anticipados a partir de la telemetría. Los registros describen eventos, las métricas resumen series numéricas y las trazas conectan el recorrido de una solicitud. Ninguna señal basta por sí sola. Un identificador de traza común permite pasar de una alerta agregada a una operación concreta.

\begin{figure}[htbp]
\centering
\begin{tikzpicture}[node distance=1.1cm and 1.1cm, every node/.style={font=\small}]
\node[system] (app) {Servicio instrumentado};
\node[activity, right=of app] (otel) {OpenTelemetry Collector};
\node[evidence, above right=0.6cm and 1.2cm of otel] (prom) {Prometheus\\métricas};
\node[evidence, right=of otel] (jaeger) {Jaeger\\trazas};
\node[evidence, below right=0.6cm and 1.2cm of otel] (logs) {Backend\\registros};
\node[decision, right=2.0cm of jaeger] (graf) {Grafana\\paneles y alertas};
\draw[arrow] (app)--node[above]{OTLP}(otel);
\draw[arrow] (otel)--(prom); \draw[arrow] (otel)--(jaeger); \draw[arrow] (otel)--(logs);
\draw[arrow] (prom)--(graf); \draw[arrow] (jaeger)--(graf); \draw[arrow] (logs)--(graf);
\end{tikzpicture}
\caption{Flujo conceptual de telemetría; los backends concretos son sustituibles.}
\end{figure}

\section{Diseño de señales}

Para servicios orientados a solicitudes, RED propone tasa, errores y duración. Para recursos, USE propone utilización, saturación y errores. Una métrica debe tener nombre, unidad, descripción y cardinalidad controlada. Etiquetar con identificador de usuario o URL completa crea series potencialmente ilimitadas y puede exponer datos personales.

Los registros estructurados facilitan consulta. Deben incluir marca temporal, nivel, evento, componente e identificadores de correlación, pero no tokens, contraseñas, cuerpos completos ni datos personales innecesarios. Las trazas se muestrean: conservar el 100\% puede ser inviable; el muestreo debe mantener suficientes errores y casos lentos para investigar.

\section{SLI, SLO y presupuesto de error}

Un indicador de nivel de servicio (SLI) cuantifica una experiencia: proporción de solicitudes válidas atendidas correctamente antes de 400 ms. Un objetivo (SLO) fija el nivel esperado en una ventana: 99,5\% durante 28 días. El acuerdo (SLA) añade consecuencias contractuales y no debe confundirse con el objetivo interno.

Si el SLO mensual es 99,5\%, el presupuesto de error es 0,5\%. Para una ventana de 28 días, esto equivale a:
\[
28\times24\times60\times0{,}005=201{,}6\text{ minutos}.
\]
El cálculo no autoriza consumir indiscriminadamente esa indisponibilidad; sirve para equilibrar confiabilidad y velocidad de cambio.

\begin{table}[htbp]
\centering
\caption{Especificación de un SLI útil.}
\begin{tabularx}{\textwidth}{p{3.2cm}X}
\toprule
\textbf{Elemento} & \textbf{Definición para LibreReserva}\\
\midrule
Evento válido & Solicitud HTTP a crear o consultar reserva, excluyendo tráfico sintético identificado.\\
Evento bueno & Respuesta no 5xx completada en menos de 400 ms.\\
SLI & Eventos buenos divididos por eventos válidos.\\
SLO & Al menos 99,5\% en ventana móvil de 28 días.\\
Fuente & Métricas de servidor instrumentadas con OpenTelemetry.\\
Respuesta & Congelar cambios riesgosos si la tasa de consumo amenaza agotar el presupuesto.\\
\bottomrule
\end{tabularx}
\end{table}

\section{Instrumentación libre y neutral}

OpenTelemetry define APIs, SDK, convenciones y un protocolo para producir y transportar telemetría sin atar el código a un proveedor. El Collector recibe, procesa y exporta señales. Prometheus consulta y almacena métricas; Grafana visualiza; Jaeger permite explorar trazas. La demostración oficial de OpenTelemetry ofrece un sistema distribuido deliberadamente complejo para experimentar.

\begin{instalacion}
Para la aplicación Python:
\begin{lstlisting}[language=bash]
python -m pip install opentelemetry-distro \
  opentelemetry-exporter-otlp
opentelemetry-bootstrap -a install
\end{lstlisting}
Para un laboratorio integrado use Podman Compose o Docker Compose y los archivos de la demostración oficial: \url{https://opentelemetry.io/docs/demo/}. Prometheus documenta sus imágenes en \url{https://prometheus.io/docs/prometheus/latest/installation/}; Grafana, en \url{https://grafana.com/docs/grafana/latest/setup-grafana/configure-docker/}. Compruebe puertos disponibles y memoria: el demo completo puede exceder 6 GiB; para equipos limitados despliegue solo aplicación, Collector, Prometheus y Jaeger.
\end{instalacion}

\begin{lstlisting}[language=bash,caption={Ejecución de FastAPI con instrumentación automática.}]
export OTEL_SERVICE_NAME=librereserva-api
export OTEL_EXPORTER_OTLP_ENDPOINT=http://localhost:4317
export OTEL_EXPORTER_OTLP_PROTOCOL=grpc
opentelemetry-instrument uvicorn app.main:app --port 8000
\end{lstlisting}

\section{Validar con una hipótesis operativa}

Un panel decorativo no constituye observabilidad. Cada panel debe apoyar una decisión. Ejemplo: ``el aumento del p95 se origina en esperas de base de datos''. La métrica identifica la ventana; una traza lenta muestra el tramo; el registro estructurado confirma el evento; una prueba controlada valida la causa.

\begin{validacion}
Genere diez solicitudes correctas, una que produzca conflicto 409 y, en una rama experimental, una latencia artificial. Verifique: (1) la métrica de tasa aumenta; (2) 409 no se clasifica automáticamente como fallo del servidor; (3) la traza lenta contiene el tramo responsable; (4) el registro comparte \texttt{trace\_id}; (5) ningún atributo contiene secretos o datos personales.
\end{validacion}

\section{Rendimiento, resiliencia y costo}

La latencia se informa con percentiles, no solo promedio. Un límite de tiempo evita esperas infinitas; un reintento solo es seguro cuando la operación lo permite, usa retroceso y dispersión, y tiene un presupuesto. De lo contrario amplifica una sobrecarga. El circuito de corte, el aislamiento de recursos y la degradación controlada reducen propagación de fallos.

El costo también es atributo arquitectónico. Réplicas o trazas adicionales pueden mejorar confiabilidad y elevar consumo. Debe reportarse por unidad útil ---por ejemplo, costo estimado por mil reservas--- y relacionarse con el SLO. En laboratorio se aceptan aproximaciones basadas en CPU, memoria y almacenamiento; deben explicitar supuestos.

\begin{ejercicio}
\textbf{Laboratorio 9.1 --- pregunta antes que panel.} Instrumente LibreReserva, construya un panel RED y documente una consulta de diagnóstico que enlace métrica, traza y registro. Incluya capturas, configuración versionada y explicación; no evalúe estética aislada.

\textbf{Laboratorio 9.2 --- experimento de resiliencia.} Inyecte demora y fallos a una dependencia local, compare sin y con timeout/reintento acotado y grafique p50, p95, errores y volumen adicional. Concluya cuándo el reintento ayudó y cuándo empeoró el sistema.
\end{ejercicio}

\section{Preguntas de revisión}
\begin{enumerate}
  \item ¿Qué etiqueta de alta cardinalidad eliminaría y qué alternativa conservaría capacidad diagnóstica?
  \item ¿Cómo cambia una decisión de despliegue cuando queda 10\% del presupuesto de error?
  \item ¿Por qué el promedio de latencia puede ocultar una mala experiencia real?
\end{enumerate}
